\documentclass[11pt]{article}
\usepackage[utf8]{inputenc}
\usepackage[T1]{fontenc}
\usepackage[spanish]{babel}
\usepackage{amsmath, amssymb}
\usepackage[ruled,vlined,linesnumbered]{algorithm2e}
\usepackage[margin=2.5cm]{geometry}

\SetKwInOut{Input}{Entrada}
\SetKwInOut{Output}{Salida}
\SetKw{Continue}{continuar}
\SetKwFunction{RadiusNeighbors}{VecindariosPorRadio}
\SetKwFunction{Push}{apilar}
\SetKwFunction{Pop}{desapilar}

\title{Pseudoc\'odigo de DBSCAN y de sus b\'usquedas de vecinos\\
\large Experimento E0 --- Implementaci\'on propia (\texttt{dbscan\_implementation})}
\date{}

\begin{document}
\maketitle

\section{Definiciones}

Sea $X=\{x_1,\dots,x_n\}$, una distancia $d$, un radio $\varepsilon>0$ y un
m\'inimo $\mathit{minPts}\ge 1$ (Ester et~al., 1996). El vecindario de un punto
\emph{incluye al propio punto}:
\[
  N_\varepsilon(x_i)=\{\,x_j : d(x_i,x_j)\le\varepsilon\,\}.
\]
\begin{itemize}
  \item $x_i$ es \emph{n\'ucleo} si $|N_\varepsilon(x_i)|\ge\mathit{minPts}$.
  \item $x_j$ es \emph{frontera} si no es n\'ucleo pero pertenece al
        vecindario de alg\'un n\'ucleo.
  \item El resto de puntos son \emph{ruido} (etiqueta $-1$).
\end{itemize}

Para vectores normalizados $\|x\|=\|y\|=1$ se cumple
\[
  \|x-y\|_2^2 = 2\,(1-\cos(x,y)) = 2\,d_{\cos}(x,y),
\]
por lo que DBSCAN con $\varepsilon_2=\sqrt{2\,\varepsilon_{\cos}}$ produce la
misma partici\'on que DBSCAN con distancia coseno y $\varepsilon_{\cos}$.

\section{Algoritmos}

\begin{algorithm}[H]
\caption{DBSCAN (\texttt{DBSCAN.fit}, \texttt{\_expand\_clusters})}
\Input{$X\in\mathbb{R}^{n\times m}$, $\varepsilon$, $\mathit{minPts}$,
       distancia $d$, b\'usqueda de vecinos, \texttt{deduplicar} (booleano)}
\Output{etiquetas $\ell\in\{-1,0,1,\dots\}^n$, \'indices de los n\'ucleos}
$N \leftarrow$ \RadiusNeighbors{$X,\ d,\ \varepsilon$}\tcp*{lista de $n$ vecindarios}
\For{$i\leftarrow 1$ \KwTo $n$}{
  $\mathit{core}_i \leftarrow \big(|N_i|\ge \mathit{minPts}\big)$\;
}
$\ell_i \leftarrow -1$ para todo $i$;\quad $c \leftarrow 0$;\quad
$\mathit{encolado}_i \leftarrow \text{falso}$ para todo $i$\;
\For{$s\leftarrow 1$ \KwTo $n$ \tcp*[f]{en orden de \'indice}}{
  \lIf{$\ell_s\ne -1$ \textbf{o} no $\mathit{core}_s$}{\Continue}
  $P \leftarrow [s]$\tcp*{pila LIFO}
  $\mathit{encolado}_s \leftarrow \text{verdadero}$\;
  \While{$P$ no est\'a vac\'ia}{
    $q \leftarrow$ \Pop{$P$}\;
    \lIf{$\ell_q\ne -1$}{\Continue}
    $\ell_q \leftarrow c$\;
    \If{$\mathit{core}_q$}{
      \ForEach{$r\in N_q$}{
        \If{$\ell_r=-1$ \textbf{y} (no \texttt{deduplicar} \textbf{o} no $\mathit{encolado}_r$)}{
          \Push{$P,\ r$};\quad $\mathit{encolado}_r \leftarrow \text{verdadero}$\;
        }
      }
    }
  }
  $c \leftarrow c+1$\;
}
\Return{$\ell$,\ $\{i : \mathit{core}_i\}$}
\end{algorithm}

\noindent\textbf{Notas.} Solo los n\'ucleos ampl\'ian el cluster: una frontera
recibe la etiqueta pero no arrastra a sus vecinos. Una frontera alcanzable
desde n\'ucleos de dos clusters se queda con el primero que la visita, por eso
su etiqueta depende del orden de los datos (indeterminaci\'on conocida de
DBSCAN). Recorrer los \'indices en orden y usar una pila LIFO reproduce la
numeraci\'on de \texttt{sklearn.cluster.DBSCAN}. La opci\'on \texttt{deduplicar}
no cambia las etiquetas: solo evita apilar dos veces un mismo punto.

\begin{algorithm}[H]
\caption{Vecindarios por bloques (\texttt{BlockedBruteForceSearch})}
\Input{$X\in\mathbb{R}^{n\times m}$, distancia $d$, radio $\varepsilon$,
       presupuesto $B$ en bytes (por defecto 256 MiB)}
\Output{lista $N$ con el vecindario de cada punto}
$b \leftarrow \max\big(1,\ \lfloor B/(8\,n)\rfloor\big)$\tcp*{filas por bloque}
$N \leftarrow [\ ]$\;
\For{$s \leftarrow 1$ \KwTo $n$ \textbf{paso} $b$}{
  $D \leftarrow d\big(X_{s:s+b-1},\ X\big)$\tcp*{matriz $b\times n$}
  \ForEach{fila $D_k$ de $D$}{
    a\~nadir a $N$ el conjunto $\{\,j : D_{kj}\le\varepsilon\,\}$\;
  }
  liberar $D$\;
}
\Return{$N$}
\end{algorithm}

\noindent Memoria $O(b\,n + n\bar k)$ con $\bar k$ el tama\~no medio del
vecindario, frente a $O(n^2)$ de la b\'usqueda por fuerza bruta que guarda la
matriz completa. El tiempo sigue siendo $O(n^2 m)$.

\begin{algorithm}[H]
\caption{Vecindarios por bloques optimizados (\texttt{OptimizedBlockedBruteForceSearch})}
\Input{$X\in\mathbb{R}^{n\times m}$, distancia $d$, radio $\varepsilon$, tama\~no de bloque $b$}
\Output{lista $N$ con el vecindario de cada punto}
\uIf{$d$ es euclídea}{
  $\tau \leftarrow \varepsilon^2$\;
  \For{cada bloque $A = X_{s:s+b-1}$}{
    $D_{kj} \leftarrow \max\big(0,\ \|a_k\|^2+\|x_j\|^2-2\,a_k\!\cdot\! x_j\big)$\tcp*{sin ra\'iz cuadrada}
    $N_k \leftarrow \{\,j : D_{kj}\le\tau\,\}$\;
  }
}
\uElseIf{$d$ es coseno}{
  $\hat X \leftarrow$ filas de $X$ divididas por su norma (norma $0\to 1$), una sola vez\;
  $\tau \leftarrow 1-\varepsilon$\;
  \For{cada bloque $\hat A = \hat X_{s:s+b-1}$}{
    $S \leftarrow \hat A\,\hat X^{\top}$\tcp*{similitudes}
    $N_k \leftarrow \{\,j : S_{kj}\ge\tau\,\}$\;
  }
}
\Else{
  usar la ruta gen\'erica de la b\'usqueda por bloques (algoritmo anterior)\;
}
\Return{$N$}
\end{algorithm}

\noindent En ambos casos la comparaci\'on es equivalente a $d(x_i,x_j)\le
\varepsilon$ en aritm\'etica exacta, pero con coma flotante los pares con
$d\approx\varepsilon$ pueden clasificarse distinto que en scikit-learn
(vease el experimento E0).

\begin{algorithm}[H]
\caption{Asignaci\'on de una instancia nueva (dise\~no, a\'un no implementado)}
\Input{instancia $q$, puntos n\'ucleo $\mathcal{C}$ con etiquetas $c(\cdot)$,
       $\varepsilon$, distancia $d$, textos originales}
\Output{cluster de $q$ y los $k$ textos m\'as cercanos de ese cluster}
$p^\ast \leftarrow \arg\min_{p\in\mathcal{C}} d(q,p)$\;
\uIf{$d(q,p^\ast)\le\varepsilon$}{
  $\hat c \leftarrow c(p^\ast)$\;
  $V \leftarrow$ los $k$ puntos de $\mathcal{C}$ con etiqueta $\hat c$ m\'as cercanos a $q$\;
}
\Else{
  $\hat c \leftarrow -1$;\quad $V \leftarrow \emptyset$\;
}
\Return{$\hat c$,\ textos de $V$}
\end{algorithm}

\end{document}
